\chapter{Evaluation in Realistic Video Scenario}
\label{c:video}

\section{Motivation and Questions}
Although Chapter~\ref{c:optoflood} quantifies recovery in terms of forwarding-level continuity and bounded overhead, real-time applications ultimately care about user-perceived smoothness. In live video delivery, a short disruption can translate into visible stalls, delayed playback progress, and burst recovery. This chapter therefore evaluates OptoFlood under a realistic video workload, focusing on (i) stall behaviour and recovery, and (ii) end-to-end delivery stability across producer handoffs.

\section{Video Delivery Model over NDN}
\subsection{Name Design and Chunking}
\subsection{Consumer Scheduling and Playback Buffer}
\subsection{Producer Behaviour under Mobility}

\section{Experimental Scenario: Access Micro-mobility}
\subsection{Topology and Mobility Trace}
\subsection{Compared Schemes}

\section{Methodology and Instrumentation}

\section{Metrics: From Continuity to QoE}
\subsection{Startup and Stall Metrics}
\subsection{Playback Recovery and Stability}
\subsection{Relating QoE to Forwarding Metrics}

\section{Results and Analysis Plan}

\section{Discussion}

\section{Summary}